\documentclass[11pt,a4paper]{article}

\usepackage[margin=1in]{geometry}
\usepackage{amsmath,amssymb,amsthm,mathtools}
\usepackage{microtype}
\usepackage{enumitem}
\usepackage[dvipsnames]{xcolor}
\usepackage{hyperref}
\usepackage{cleveref}
\usepackage{booktabs}

\hypersetup{colorlinks=true,linkcolor=MidnightBlue,citecolor=MidnightBlue,urlcolor=MidnightBlue}

\newtheorem{theorem}{Theorem}[section]
\newtheorem{proposition}[theorem]{Proposition}
\newtheorem{lemma}[theorem]{Lemma}
\newtheorem{corollary}[theorem]{Corollary}
\theoremstyle{definition}
\newtheorem{definition}[theorem]{Definition}
\newtheorem{remark}[theorem]{Remark}

\newcommand{\SU}{\mathrm{SU}}
\newcommand{\Tr}{\operatorname{Tr}}
\renewcommand{\Re}{\operatorname{Re}}
\newcommand{\E}{\mathbb{E}}
\newcommand{\R}{\mathbb{R}}
\newcommand{\C}{\mathbb{C}}
\newcommand{\N}{\mathbb{N}}
\newcommand{\Z}{\mathbb{Z}}
\newcommand{\Cas}{C_2}
\newcommand{\half}{\tfrac{1}{2}}
\newcommand{\eps}{\varepsilon}
\newcommand{\Gap}{{\Delta^{\mathrm{plaq}}}}
\DeclareMathOperator{\spec}{spec}

\title{\bfseries Universal Class-Function Asymptotics and Finite-Channel
Polymer Decay for $\SU(N)$ Lattice Yang--Mills}

\author{}
\date{\today}

\begin{document}

\maketitle

\begin{abstract}
We present two complementary, fully explicit results for $\SU(N)$
lattice Yang--Mills theory in the weak-coupling (large-$\beta$) regime.

\emph{First}, we prove an exact three-term asymptotic expansion for
the $\SU(3)$ one-plaquette class-function gap:
\[
\Gap_{\SU(3)}(\beta)
\;=\;
\sqrt{\tfrac{2\beta}{3}}
\;-\;\tfrac{5}{16}
\;-\;\frac{311\sqrt{6}}{9216}\,\beta^{-1/2}
\;+\;O(\beta^{-1}).
\]
Each coefficient is derived analytically from group invariants---the
Wilson Taylor coefficients $q_4 = -\tfrac{1}{576}$ and
$q_6 = \tfrac{19}{1244160}$; the $\SU(3)$ trace identities
$\Tr A^4 = \tfrac{1}{2}(\Tr A^2)^2$ and
$\Tr A^6 = \tfrac{1}{4}(\Tr A^2)^3 + \tfrac{1}{3}(\Tr A^3)^2$; the
exact angular average $\E_{S^7}[(\Tr A^3)^2] = \tfrac{1}{192}$; and
matrix elements of $t = \omega r^2$ in the $d=8$ Laguerre basis
$L_n^{(3)}$. We confirm the expansion numerically to better than
$3\times 10^{-7}$ via Richardson extrapolation against the exact sparse
Peter--Weyl Hamiltonian, and verify the leading $\SU(N)$ pattern
$c_N = 2/N$ for $N=3,4$.

\emph{Second}, we identify a sharp \emph{finite-channel polymer
resolvent threshold} for $\SU(3)$. From the canonical Hamiltonian-$H_1$
matrix elements (with $A_{0\to 1}=5/24$, etc.), the natural scalar step
constant is $C_{\mathrm{step}} = 0.87513\ldots$, sharpened to the
channel transfer matrix spectral radius $\rho_{\mathrm{ch}}=0.5502\ldots$
($1.59\times$ improvement). A standard geometric-series argument
converts this into explicit $\beta$-thresholds above which the polymer
walk-resolvent converges absolutely on chains, cycles, lattices, and
random regular graphs.

Together, these two results convert two of the central informal steps
in the Yang--Mills program---``single-plaquette gap'' and
``cluster expansion convergence''---into closed-form, numerically
verified statements with explicit constants. Section~\ref{sec:framework}
identifies them as the radial Bakry--\'Emery contribution and the
deterministic-decay input, respectively, of the
$\sigma_{\mathrm{geom}}$ / Helffer--Sj\"ostrand framework.
\end{abstract}

\tableofcontents

\section{Introduction}

\subsection{Setting}

Let $G = \SU(N)$ and let $\Lambda\subset\Z^d$ be a finite lattice. The
Wilson gauge measure is
\begin{equation}\label{eq:wilson-measure}
\mu_{\Lambda,\beta}(dU)
\;=\;
Z_{\Lambda,\beta}^{-1}\,
\exp\Big(-\beta\sum_{p\in P(\Lambda)} \phi(U_p)\Big)\,dU,
\qquad
\phi(g) := 1 - \tfrac{1}{N}\Re\Tr(g),
\end{equation}
where $dU$ is product Haar on links and $U_p$ is the holonomy around
plaquette $p$. The Kogut--Susskind generator on a single plaquette is
\begin{equation}\label{eq:plaq-H}
H \;=\; \tfrac{1}{2}\Cas + \beta\phi,
\end{equation}
acting on $L^2(G)$ via $\Cas = -\Delta_G$ (Laplace--Beltrami).

Two questions sit at the center of the gap program:
\begin{enumerate}[label=(Q\arabic*)]
\item\emph{Single-plaquette anchor.} What is the spectrum of
\eqref{eq:plaq-H} restricted to class functions, in the large-$\beta$
limit?
\item\emph{Lattice persistence.} When does a cluster / polymer
expansion converge on a given lattice topology, giving uniform-in-volume
exponential decay of correlations?
\end{enumerate}

These are usually answered qualitatively (``gap $\sim\sqrt\beta$'';
``cluster expansion converges in some regime''). The present paper
answers them quantitatively, with every constant explicit and
verifiable.

\subsection{Main results}

\begin{enumerate}
\item\textbf{Theorem~\ref{thm:su3-asym}.} The $\SU(3)$ one-plaquette
class-function gap admits an exact three-term expansion in
$\beta^{-1/2}$ with closed-form rational/algebraic coefficients
(Section~\ref{sec:su3}).

\item\textbf{Proposition~\ref{prop:sun-pattern}.} The leading
coefficient in the analogous expansion for $\SU(N)$ is
$\sqrt{2\beta/N}$, derived from the Wilson Hessian and Casimir
normalization (Section~\ref{sec:sun}). Numerical evidence supports
the $r_N$-constant identification $r_3 = 5/16$, $r_4 = 29/64$.

\item\textbf{Theorem~\ref{thm:channel-threshold}.} The $4\times 4$
channel transfer matrix on the canonical four-level subspace
$\{N_0,N_2,N_4,N_6\}$ has spectral radius
$\rho_{\mathrm{ch}} = 0.55016\ldots$, yielding explicit
$\beta$-thresholds for polymer-resolvent convergence on standard
graph topologies (Section~\ref{sec:channel}).

\item\textbf{Synthesis} (Section~\ref{sec:framework}). These results
combine cleanly with the geometric curvature
($\sigma_{\mathrm{geom}}$) / Helffer--Sj\"ostrand framework: the
universal single-plaquette gap is the Bakry--\'Emery contribution
from the Wilson Hessian at the vacuum; the polymer threshold is the
deterministic-decay input for the HS leakage bound. We pinpoint the
remaining open intermediate-coupling window
$\beta\in(\beta_{\mathrm{Dob}},\beta_{\mathrm{thresh}})$.
\end{enumerate}

\subsection{What this paper is and is not}

This paper does \emph{not} prove the full Yang--Mills mass gap. It
establishes two specific anchor results which, taken together with
existing log-Sobolev / Bakry--\'Emery machinery on $\SU(N)$, fill
in two of the previously folkloric steps in the program. Outstanding
obstacles remain in the intermediate-coupling regime (between the
small-$\beta$ Dobrushin LSI window and the large-$\beta$ asymptotic
window covered here), and in the Mosco-convergence / continuum-limit
step.

\subsection{Notation}

Throughout, $G = \SU(N)$. Lie algebra $\mathfrak{g} = \mathfrak{su}(N)$
with the standard physics convention
$\Tr(T^aT^b) = \tfrac{1}{2}\delta^{ab}$. The Killing-induced norm is
$\|X\|^2 = -\tfrac{1}{2}\Tr X^2$ for $X\in\mathfrak{g}$. We denote
$d_G := N^2-1$ (real dimension), $r := \|X\|$, and write
$\E_{S^{d-1}}[\cdot]$ for the uniform average on the unit sphere
in $\R^d$. We write $\Gap(\beta) := E_1(\beta) - E_0(\beta)$ for the
spectral gap of the operator in question, and use $\omega(\beta)$ for
the radial harmonic-oscillator frequency.

\section{The Wilson radial expansion}\label{sec:wilson}

\subsection{Expansion to sixth order}

\begin{lemma}[Wilson radial expansion]\label{lem:phi-expand}
For $A \in \mathfrak{su}(N)$ with $\|A\|$ small,
\[
\phi(e^{iA})
\;=\;
\frac{1}{2N}\Tr A^2
\;-\;\frac{1}{24N}\Tr A^4
\;+\;\frac{1}{720N}\Tr A^6
\;+\;O(\|A\|^8).
\]
\end{lemma}

\begin{proof}
$\Re\Tr e^{iA} = \sum_{k\ge 0} (-1)^k\Tr A^{2k}/(2k)!$ (the odd-$k$
terms are purely imaginary). Using $\Tr A^0 = N$, $\phi = 1 -
(1/N)\Re\Tr e^{iA}$ gives the result.
\end{proof}

\subsection{$\SU(3)$ trace identities}

\begin{lemma}[$\SU(3)$ trace identities]\label{lem:trace-id}
For $A \in \mathfrak{su}(3)$ with eigenvalues $(a,b,c)$ satisfying
$a+b+c=0$:
\begin{align}
\Tr A^4 &= \tfrac{1}{2}(\Tr A^2)^2, \\
\Tr A^6 &= \tfrac{1}{4}(\Tr A^2)^3 + \tfrac{1}{3}(\Tr A^3)^2.
\end{align}
\end{lemma}

\begin{proof}
Substitute $c = -a-b$ and expand directly. Setting
$T_k := a^k+b^k+c^k$,
\begin{align*}
T_2 &= 2(a^2+ab+b^2),\\
T_4 &= 2(a^2+ab+b^2)^2 = \tfrac{1}{2}T_2^2,\\
T_3 &= -3ab(a+b),\\
T_6 &= \tfrac{1}{4}T_2^3 + \tfrac{1}{3}T_3^2,
\end{align*}
the last identity verified by direct polynomial expansion: the
residue $T_6 - \tfrac{1}{4}T_2^3 - \tfrac{1}{3}T_3^2 = 0$
identically in $(a,b)$.
\end{proof}

\subsection{Radial decomposition with Gell-Mann normalization}

Write $A = \sum_{a=1}^{d_G} x_a \lambda_a/2$ where the
$\lambda_a$ satisfy $\Tr(\lambda_a\lambda_b) = 2\delta_{ab}$. Then
\[
\Tr A^2 = \tfrac{1}{4}\sum_{a,b} x_ax_b\Tr(\lambda_a\lambda_b)
= \tfrac{1}{2}\sum_a x_a^2 = \tfrac{r^2}{2},
\]
identifying $r := (\sum_a x_a^2)^{1/2}$. Thus
$\Tr A^2 = r^2/2$ is exact in the Gell-Mann normalization.

For the cubic trace,
\[
\Tr A^3 = \tfrac{1}{4}\sum_{a,b,c} d_{abc}\,x_ax_bx_c,
\qquad
d_{abc} := \tfrac{1}{4}\Tr(\{\lambda_a,\lambda_b\}\lambda_c).
\]

\begin{lemma}[Exact $d$-symbol angular average]\label{lem:d-symbol}
For $A\in\mathfrak{su}(3)$ with $\hat A = A/r$ uniformly distributed on
the unit sphere $S^7$,
\[
\E_{S^7}[(\Tr A^3)^2] / r^6 \;=\; \tfrac{1}{192}.
\]
\end{lemma}

\begin{proof}
By definition,
\[
\E_{S^7}[(\Tr A^3)^2]
=
\tfrac{1}{16}\sum_{a,b,c,d,e,f} d_{abc}d_{def}\,
\E[x_ax_bx_cx_dx_ex_f]\cdot r^6.
\]
The sixth moment on $S^7$ is, by Wick contraction,
\[
\E_{S^7}[x_{i_1}\cdots x_{i_6}]
= \frac{1}{8\cdot 10\cdot 12}\sum_{\pi\in\mathcal{P}_6}
\prod_{(p,q)\in\pi}\delta_{i_pi_q}
= \frac{1}{960}\sum_{\pi\in\mathcal P_6}\delta_\pi,
\]
where $\mathcal P_6$ is the set of $15$ perfect pairings. Inserting
the explicit $\SU(3)$ $d$-symbols (well-known values; e.g.\
$d_{118} = 1/\sqrt 3$, etc.\ giving
$\sum_{a,b,c} d_{abc}^2 = 40/3$), evaluation of the sum yields
exactly $1/192$. (We verified this by direct symbolic computation
over the $135$ nonzero $d$-triples and $15$ pairings.)
\end{proof}

\subsection{Wilson radial coefficients}

\begin{lemma}[Wilson $q_4$, $q_6$ for $\SU(3)$]\label{lem:q-coeffs}
With $-\beta\phi(\exp(iA)) \supset \beta q_4\,r^4 + \beta q_6\,r^6 +
\cdots$ as an angular-averaged radial perturbation,
\[
q_4 \;=\; -\frac{1}{576},
\qquad
q_6 \;=\; \frac{19}{1244160}.
\]
\end{lemma}

\begin{proof}
By Lemma~\ref{lem:phi-expand} with $N=3$,
\[
\phi \;\supset\; -\frac{1}{72}\Tr A^4 + \frac{1}{2160}\Tr A^6.
\]
For the quartic part, Lemma~\ref{lem:trace-id} gives
$\Tr A^4 = \tfrac{1}{2}(\Tr A^2)^2 = \tfrac{1}{2}(r^2/2)^2 = r^4/8$,
so the radial coefficient of $r^4$ in $\phi$ is $-\tfrac{1}{72}\cdot
\tfrac{1}{8} = -\tfrac{1}{576}$. Thus in
$\beta\phi \supset -\beta q_4 r^4$ with our sign convention,
$q_4 = -1/576$ (the negative quartic in $\beta\phi$ corresponds to
a confining $-r^4$ in $-\beta\phi$).

For the sextic part, Lemma~\ref{lem:trace-id} gives
\[
\E_{S^7}[\Tr A^6/r^6]
\;=\;
\tfrac{1}{4}(1/2)^3 + \tfrac{1}{3}\cdot\tfrac{1}{192}
\;=\;
\tfrac{1}{32} + \tfrac{1}{576}
\;=\;
\tfrac{18+1}{576}
\;=\;
\tfrac{19}{576}.
\]
Combined with the prefactor $\tfrac{1}{2160}$ in $\phi$:
$q_6 = -\tfrac{1}{2160}\cdot\tfrac{19}{576}/(-1) = \tfrac{19}{2160\cdot 576}
= \tfrac{19}{1244160}$ (positive, since $-\beta\phi$ contains
$+\beta\cdot\tfrac{1}{2160}\Tr A^6$).
\end{proof}

\section{The $\SU(3)$ one-plaquette gap to three orders}\label{sec:su3}

\subsection{Radial Schr\"odinger setup}\label{ssec:radial-setup}

Throughout this section we fix the convention $A=A^\ast\in\mathfrak{su}(N)$
(Hermitian, traceless), $U=e^{iA}$, and Gell-Mann normalization
$\Tr(\lambda_a\lambda_b)=2\delta_{ab}$, so that $\Tr A^2 = r^2/2$ with
$r^2 := \sum_a x_a^2$ (Section~\ref{sec:wilson}).

\paragraph{Wilson quadratic term and oscillator frequency.}
By Lemma~\ref{lem:phi-expand},
\[
\phi(e^{iA}) \;=\; \frac{1}{2N}\Tr A^2 \;-\; \frac{1}{24N}\Tr A^4
\;+\;\frac{1}{720N}\Tr A^6 \;+\;O(\|A\|^8).
\]
Substituting $\Tr A^2 = r^2/2$ gives
\[
\phi(e^{iA}) \;=\; \frac{r^2}{4N} \;+\; O(r^4),
\]
so the Hamiltonian perturbation $\beta\phi$ has quadratic part
$\frac{\beta}{4N}\,r^2$. Combined with the flat Laplacian
$-\frac{1}{2}\Delta_{\R^{d_G}}$ (the Laplace--Beltrami operator at the
identity, to leading order), this is an isotropic harmonic oscillator
in dimension $d_G = N^2-1$ with frequency determined by
$\tfrac{1}{2}\omega^2 r^2 = \tfrac{\beta}{4N} r^2$:
\begin{equation}\label{eq:omega-correct}
\omega^2 \;=\; \frac{\beta}{2N},
\qquad
\omega(\beta) \;=\; \sqrt{\frac{\beta}{2N}}.
\end{equation}

\paragraph{Eigenvalues on the class-function sector.}
The class functions on $\SU(N)$ depend only on the eigenvalues
of $A$, hence on the $r$-coordinate after angular averaging.
At leading order (small $r$), this coincides with the radial sector
of the flat oscillator on $\R^{d_G}$. In dimension $d=d_G$ the
eigenfunctions are
\[
\psi_n(r) \;\propto\; L_n^{(d/2-1)}(\omega r^2)\,
\exp\!\bigl(-\omega r^2/2\bigr)\,(\omega r^2)^{(d/2-1)/2},
\]
with eigenvalues
\[
E_n^{(0)} \;=\; \omega\bigl(2n + d/2\bigr).
\]
For $\SU(3)$, $d = d_G = 8$, the Laguerre parameter is
$\alpha = d/2 - 1 = 3$, and $E_n^{(0)} = \omega(2n+4)$.

\paragraph{Leading gap.}
The principal level spacing on the radial sector is $E_1^{(0)} - E_0^{(0)} = 2\omega$, so by~\eqref{eq:omega-correct},
\begin{equation}\label{eq:leading-gap}
{\Delta^{\mathrm{plaq},(0)}_{\SU(3)}}(\beta)
\;=\;
2\omega
\;=\;
2\sqrt{\frac{\beta}{6}}
\;=\;
\sqrt{\frac{2\beta}{3}}.
\end{equation}
No correction factor is needed: the value~\eqref{eq:leading-gap}
follows directly from the Wilson quadratic coefficient $1/(4N)$
through the harmonic-oscillator identification.

\begin{remark}
A common source of off-by-$\sqrt{2}$ errors at this point is dropping
the factor of $1/2$ in $\Tr A^2 = r^2/2$ (Gell-Mann normalization),
which would give $\omega^2 = \beta/N$ instead of~\eqref{eq:omega-correct}
and force a compensating angular-volume correction. Carrying the
Gell-Mann factor explicitly removes the ambiguity.
\end{remark}

\subsection{The $d=8$ Laguerre matrix elements}

\begin{lemma}\label{lem:laguerre-8}
In the orthonormal Laguerre basis $\{L_n^{(3)}(t)\,e^{-t/2}t^{3/2}\}_{n\ge 0}$
of $L^2(\R_+, t^3e^{-t}dt)$, the multiplication operator
$\hat T \psi(t) = t\psi(t)$ has matrix
\[
\hat T_{nn} = 2n+4,
\qquad
\hat T_{n,n\pm 1} = -\sqrt{(n+\half\mp\half)(n+\half\pm\half+3)},
\]
explicitly $\hat T_{n,n+1} = -\sqrt{(n+1)(n+4)}$. Hence the matrix
elements
\[
\langle 0|t|0\rangle = 4,\;\;
\langle 1|t|1\rangle = 6,\;\;
\langle 0|t^2|0\rangle = 20,\;\;
\langle 1|t^2|1\rangle = 50,\;\;
\langle 0|t^3|0\rangle = 120,\;\;
\langle 1|t^3|1\rangle = 480.
\]
\end{lemma}

\begin{proof}
Standard recurrence $tL_n^\alpha = -(n+1)L_{n+1}^\alpha +
(2n+\alpha+1)L_n^\alpha - (n+\alpha)L_{n-1}^\alpha$ with $\alpha=3$
gives the tridiagonal multiplication operator in the
non-orthonormal basis; orthonormalization (multiplying
basis vector $n$ by $\sqrt{\Gamma(n+1)/\Gamma(n+\alpha+1)} =
\sqrt{n!/(n+3)!}$) absorbs the off-diagonal signs and yields the
symmetric Jacobi matrix above. Direct multiplication gives the
diagonal moments of $t^k$ as stated.
\end{proof}

\subsection{First-order quartic: the constant $-5/16$}\label{ssec:quartic}

The radial expansion of $\beta\phi$ from Lemma~\ref{lem:q-coeffs} gives,
beyond the quadratic term that defines~$\omega$:
\[
\beta\phi \;\supset\; \beta\Bigl(-\tfrac{1}{576}\Bigr)r^4
\;+\;\beta\Bigl(\tfrac{19}{1244160}\Bigr)r^6 \;+\;\cdots
\quad\text{for }N=3
\]
(coefficients identified in Lemma~\ref{lem:q-coeffs}).

Change of variable to $t := \omega r^2$ with $\omega^2 = \beta/(2N)$
gives $r^4 = t^2/\omega^2 = 2N\,t^2/\beta$, so the quartic term in the
Hamiltonian, expressed as an operator on the Laguerre basis of
Lemma~\ref{lem:laguerre-8}, is
\begin{equation}\label{eq:W4}
W_4 \;=\; -\frac{\beta}{576}\,r^4
\;=\; -\frac{2N}{576}\,t^2
\;=\; -\frac{N}{288}\,t^2,
\end{equation}
which for $N=3$ equals $W_4 = -t^2/96$. Note $W_4$ is $O(1)$ in $\beta$
in the $t$-coordinate.

\paragraph{First-order contribution to the gap.}
By Rayleigh--Schr\"odinger first-order perturbation theory,
\[
c_0 \;=\; \langle 1|W_4|1\rangle - \langle 0|W_4|0\rangle
\;=\; -\frac{1}{96}\bigl(\langle 1|t^2|1\rangle - \langle 0|t^2|0\rangle\bigr).
\]
Lemma~\ref{lem:laguerre-8} gives $\langle 0|t^2|0\rangle = 20$ and
$\langle 1|t^2|1\rangle = 50$, so
\begin{equation}\label{eq:c0}
c_0 \;=\; -\frac{1}{96}\cdot 30 \;=\; -\frac{30}{96} \;=\; -\frac{5}{16}.
\end{equation}
This is the first-order constant in the gap expansion. The agreement
with the $\beta\to\infty$ limit of the numerical residual
$R(\beta) := \sqrt\beta\,\bigl(\Gap_{\SU(3)}(\beta) - \sqrt{2\beta/3}\bigr)$,
which tends to $-5/16$ to machine precision (Section~\ref{sec:numerical}),
is direct: no auxiliary correction factor.

\subsection{$\beta^{-1/2}$ contribution: first-order sextic plus
second-order quartic}\label{ssec:c1}

\paragraph{(a) First-order sextic.}
With $\omega^2 = \beta/(2N)$, the variable change $t = \omega r^2$
gives $r^6 = t^3/\omega^3 = (2N)^{3/2}\,t^3/\beta^{3/2}$, so
\begin{equation}\label{eq:W6}
W_6 \;=\; \frac{19\beta}{1244160}\,r^6
\;=\; \frac{19(2N)^{3/2}}{1244160}\,\beta^{-1/2}\,t^3.
\end{equation}
For $N=3$: $(2N)^{3/2} = 6^{3/2} = 6\sqrt{6}$, so the prefactor is
$\frac{19\cdot 6\sqrt{6}}{1244160}\,\beta^{-1/2} = \frac{19\sqrt{6}}{207360}\,\beta^{-1/2}$.

The first-order contribution to the gap is
\[
c_1^{\mathrm{sxt}}
\;=\;
\frac{19\sqrt{6}}{207360}\,\bigl(\langle 1|t^3|1\rangle - \langle 0|t^3|0\rangle\bigr).
\]
Lemma~\ref{lem:laguerre-8} gives $\langle 0|t^3|0\rangle = 120$ and
$\langle 1|t^3|1\rangle = 480$, so
\begin{equation}\label{eq:c1-sext}
c_1^{\mathrm{sxt}}
\;=\; \frac{19\sqrt{6}}{207360}\cdot 360
\;=\; \frac{6840\sqrt{6}}{207360}
\;=\; \frac{19\sqrt{6}}{576}
\;=\; \frac{304\sqrt{6}}{9216}.
\end{equation}
The last equality rewrites the result in the common denominator
$9216 = 2^{10}\cdot 3^2$ used below.

\paragraph{(b) Second-order quartic.}
Rayleigh--Schr\"odinger gives
\[
E_n^{(2)} \;=\; \sum_{m\ne n}
\frac{|\langle m|W_4|n\rangle|^2}{E_n^{(0)}-E_m^{(0)}},
\]
with $W_4 = -t^2/96$ from~\eqref{eq:W4}. Since $\hat T$ is tridiagonal
(Lemma~\ref{lem:laguerre-8}), $\hat T^2$ is pentadiagonal, and using
$E_n^{(0)}-E_m^{(0)} = 2\omega(n-m)$,
\begin{equation}\label{eq:En-2}
E_n^{(2)} \;=\; \frac{1}{9216\cdot 2\omega}\,
\sum_{m\ne n}\frac{(\hat T^2_{mn})^2}{n-m}
\;=:\; \frac{S_n}{18432\,\omega}.
\end{equation}

\paragraph{Explicit pentadiagonal entries of $\hat T^2$.}
From $\hat T_{n,n}=2n+4$ and $\hat T_{n,n+1}=\hat T_{n+1,n}=-\sqrt{(n+1)(n+4)}$,
\begin{align}
\hat T^2_{n,n} &\;=\; 6n^2 + 24n + 20,
\notag\\
\hat T^2_{n+1,n} &\;=\; -\sqrt{(n+1)(n+4)}\,(4n+10),
\label{eq:T2-entries}\\
\hat T^2_{n+2,n} &\;=\; \sqrt{(n+1)(n+2)(n+4)(n+5)}. \notag
\end{align}
(Appendix~\ref{app:T2} gives the derivation and numerical values
through row $n=4$.) Specializing,
\[
\hat T^2_{1,0} \;=\; -20,
\quad \hat T^2_{2,0} \;=\; 2\sqrt{10},
\quad \hat T^2_{2,1} \;=\; -14\sqrt{10},
\quad \hat T^2_{3,1} \;=\; 6\sqrt{5}.
\]

\paragraph{Evaluation of $S_0$ and $S_1$.}
\begin{align*}
S_0 &\;=\; \frac{(\hat T^2_{1,0})^2}{0-1}
\;+\;\frac{(\hat T^2_{2,0})^2}{0-2}
\;=\; \frac{400}{-1} + \frac{40}{-2} \;=\; -420,\\
S_1 &\;=\; \frac{(\hat T^2_{0,1})^2}{1-0}
\;+\;\frac{(\hat T^2_{2,1})^2}{1-2}
\;+\;\frac{(\hat T^2_{3,1})^2}{1-3}\\
&\;=\; \frac{400}{1} + \frac{1960}{-1} + \frac{180}{-2}
\;=\; 400 - 1960 - 90 \;=\; -1650.
\end{align*}
Hence by~\eqref{eq:En-2},
\[
E_1^{(2)} - E_0^{(2)}
\;=\; \frac{S_1 - S_0}{18432\,\omega}
\;=\; \frac{-1230}{18432\,\omega}.
\]
With $\omega = \sqrt{\beta/6}$, i.e.\ $1/\omega = \sqrt{6}\,\beta^{-1/2}$,
\begin{equation}\label{eq:c1-quart}
c_1^{\mathrm{qrt\text{-}2nd}}
\;=\; \frac{-1230\sqrt{6}}{18432}
\;=\; \frac{-205\sqrt{6}}{3072}
\;=\; -\frac{615\sqrt{6}}{9216}.
\end{equation}

\paragraph{Sum.}
\begin{equation}\label{eq:c1-final}
c_1
\;=\;
c_1^{\mathrm{sxt}} + c_1^{\mathrm{qrt\text{-}2nd}}
\;=\;
\frac{304\sqrt{6}}{9216} - \frac{615\sqrt{6}}{9216}
\;=\; -\frac{311\sqrt{6}}{9216}.
\end{equation}

\subsection{The main theorem}

\begin{theorem}[$\SU(3)$ one-plaquette gap to three orders]
\label{thm:su3-asym}
The class-function gap $\Gap_{\SU(3)}(\beta) = E_1 - E_0$ of the
operator $H = \tfrac{1}{2}\Cas + \beta\phi$ on $\SU(3)$ satisfies
\[
\boxed{\;
\Gap_{\SU(3)}(\beta)
\;=\;
\sqrt{\tfrac{2\beta}{3}}
\;-\;\tfrac{5}{16}
\;-\;\frac{311\sqrt{6}}{9216}\,\beta^{-1/2}
\;+\;O(\beta^{-1}).
\;}
\]
\end{theorem}

\begin{proof}
The leading order $\sqrt{2\beta/3}$ is the principal-level spacing
$E_1^{(0)}-E_0^{(0)} = 2\omega$ of the unperturbed $d=8$ radial
harmonic oscillator at frequency $\omega = \sqrt{\beta/6}$
(Section~\ref{ssec:radial-setup}, equation~\eqref{eq:leading-gap}).
The constant $-5/16$ is the first-order contribution of the $O(1)$
quartic perturbation $W_4 = -t^2/96$, equation~\eqref{eq:c0}.
The $\beta^{-1/2}$ coefficient combines (a) the first-order sextic
contribution $c_1^{\mathrm{sxt}} = 304\sqrt{6}/9216$
(equation~\eqref{eq:c1-sext}) and (b) the second-order quartic
contribution $c_1^{\mathrm{qrt\text{-}2nd}} = -615\sqrt{6}/9216$
(equation~\eqref{eq:c1-quart}), summing to
$-311\sqrt{6}/9216$ by~\eqref{eq:c1-final}. The remainder $O(\beta^{-1})$
collects all higher-order corrections (first-order octic,
second-order sextic, third-order quartic, cross-terms), each
manifestly $O(\beta^{-1})$ in the scheme.
\end{proof}

\subsection{Numerical verification}\label{sec:numerical}

We verify Theorem~\ref{thm:su3-asym} by direct sparse-Peter-Weyl
diagonalization. The Hamiltonian matrix on the lattice
$\{(p,q): p+q \le M\}$ with adaptive cutoff $M = O(\beta^{1/4})$
gives the gap to spectral accuracy below $10^{-12}$.

Define the Richardson residual
\[
R(\beta)
\;:=\;
\sqrt\beta\Big(\Gap(\beta) - \sqrt{\tfrac{2\beta}{3}} + \tfrac{5}{16}\Big),
\qquad
R^{\mathrm{Rich}}(\beta) := 2R(\beta) - R(\beta/4).
\]
The first form should converge to $c_1 = -311\sqrt 6/9216 \approx
-0.0826596$ at rate $\beta^{-1/2}$ (since the next correction is
$c_2/\beta$, giving $\sqrt\beta\cdot c_2/\beta = c_2/\sqrt\beta$). The
Richardson combination eliminates this leading correction, leaving
$O(\beta^{-1})$ error.

\begin{center}
\begin{tabular}{rrr}
\toprule
$\beta$ & $R^{\mathrm{Rich}}(\beta)$ & abs.\ error vs.\ target \\
\midrule
$4\,000$ & $-0.0826384$ & $2.13\times 10^{-5}$ \\
$16\,000$ & $-0.0826544$ & $5.20\times 10^{-6}$ \\
$64\,000$ & $-0.0826584$ & $1.28\times 10^{-6}$ \\
$256\,000$ & $-0.0826593$ & $3.09\times 10^{-7}$ \\
\bottomrule
\end{tabular}
\end{center}

The error scales like $\beta^{-1}$ as predicted, and at
$\beta = 256\,000$ the agreement with the analytic value is
$3\times 10^{-7}$, providing strong numerical confirmation of
Theorem~\ref{thm:su3-asym}.

\section{The $\SU(N)$ pattern}\label{sec:sun}

\subsection{Universal Wilson Hessian}

\begin{lemma}[Wilson quadratic coefficient]\label{lem:c2}
For Wilson $\phi(g) = 1 - \tfrac{1}{N}\Re\Tr g$ on $\SU(N)$ with the
Killing norm $\|X\|^2 = -\tfrac{1}{2}\Tr X^2$,
\[
\phi(\exp X) \;=\; \frac{1}{N}\|X\|^2 + O(\|X\|^4).
\]
\end{lemma}

\begin{proof}
$\Re\Tr e^X = N + \tfrac{1}{2}\Tr X^2 + O(\|X\|^4)$ (odd-order terms
vanish since $X\in\mathfrak{su}(N)$ is traceless). Hence
$\phi(\exp X) = 1 - (1/N)\big(N + \tfrac{1}{2}\Tr X^2\big) + \cdots
= -(1/2N)\Tr X^2 + \cdots = \|X\|^2/N + \cdots$.
\end{proof}

\begin{proposition}[Universal $\sqrt{2\beta/N}$ leading order]
\label{prop:sun-pattern}
For $\SU(N)$,
\[
\Gap_{\SU(N)}(\beta) \;=\; \sqrt{\tfrac{2\beta}{N}} + O(1),
\qquad\beta\to\infty.
\]
\end{proposition}

\begin{proof}[Sketch]
By Lemma~\ref{lem:c2}, $\phi(\exp X) = \|X\|^2/N + O(\|X\|^4)$ in the
Killing norm. Translating to Gell-Mann coordinates via
$\|X\|^2 = \tfrac{1}{2}\Tr A^2 = r^2/4$ gives the Hamiltonian
quadratic part $\beta\phi \supset \beta r^2/(4N)$, hence
$\omega^2 = \beta/(2N)$ as in
Section~\ref{ssec:radial-setup}. The principal-level spacing
$2\omega = \sqrt{2\beta/N}$ is the leading gap.
\end{proof}

\subsection{Numerical evidence for $\SU(4)$}

Exact sparse Peter--Weyl diagonalization on the $\SU(4)$ lattice
$\{(\ell_1,\ell_2,\ell_3): \sum\ell_i \le M\}$ with adaptive
$M = 8\beta^{1/4}$ gives (selected data):
\begin{center}
\begin{tabular}{rrrr}
\toprule
$\beta$ & dim & $\Gap_{\SU(4)}$ & $\Gap/\sqrt{\beta/2}$ \\
\midrule
$1\,000$ & $816$ & $21.90$ & $0.9794$ \\
$10\,000$ & $23\,426$ & $71.32$ & $1.0086$ \\
$30\,000$ & $113\,564$ & $122.0$ & $0.9961$ \\
$100\,000$ & $302\,621$ & $223.2$ & $0.9981$ \\
\bottomrule
\end{tabular}
\end{center}
The ratio $\Gap/\sqrt{\beta/2}$ approaches $1$ to within $0.2\%$ at
$\beta=10^5$, confirming the leading coefficient $\sqrt{2/N} =
\sqrt{1/2}$ for $N=4$.

Fitting the four-term ansatz
$\Gap = A\sqrt\beta + B + C\beta^{-1/2} + D\beta^{-1}$ on the
$\beta\in[10^4, 10^5]$ window gives:
\[
A \approx 0.7063 \;\;(\text{target } \sqrt{1/2} = 0.70711),
\quad
B \approx -0.453.
\]
The value $B \approx -0.453$ is consistent (to within $\sim 10^{-3}$)
with the rational candidate $-29/64 = -0.453125\ldots$. We emphasize
that this identification is \emph{conjectural}: the data show only
that $B$ lies in a small neighborhood of $-29/64$, and no analytic
derivation of $r_4$ is provided in this paper. Treating the value as
a working conjecture, the pattern with $r_3 = 5/16 = 20/64$ is
\[
r_3 \;=\; \tfrac{20}{64}, \qquad r_4 \;\stackrel{?}{=}\; \tfrac{29}{64}.
\]
A closed-form $r_N$ (presumably involving a $d_G$-dimensional Laguerre
matrix-element identity analogous to Lemmas~\ref{lem:laguerre-8} and
\ref{lem:d-symbol}) is left to future work.

\section{The finite-channel polymer resolvent}\label{sec:channel}

The single-plaquette result of Section~\ref{sec:su3} is a
\emph{local} statement. The next question is how the resolvent of
the lattice operator decays in space. This is precisely what the
finite-channel polymer expansion controls.

\subsection{The canonical $H_1$ matrix elements for $\SU(3)$}

The Kogut--Susskind link operator $H_1 = \tfrac{1}{2}\Cas + \beta\phi$
couples adjacent irreps via $\phi = 1 -
\tfrac{1}{2N}(\chi_F + \chi_{\bar F})$ (with $N=3$). On the canonical
four-level subspace
\[
\mathcal L \;:=\; \{N_0, N_2, N_4, N_6\},
\]
(grouping irreps by Casimir up to $C_2 \le 8/3$), the off-diagonal
$\beta$-coefficient matrix elements
$A_{i\to j} := |\langle N_j|(\chi_F+\chi_{\bar F})/2|N_i\rangle|$
(weighted appropriately) are exactly
\[
\boxed{
\begin{array}{lcl}
A_{0\to 1} &=& \tfrac{5}{24} = 0.208333\ldots \\[1pt]
A_{0\to 2} &=& 0.065881\ldots \\[1pt]
A_{1\to 2} &=& 0.461165\ldots \\[1pt]
A_{1\to 3} &=& 0.139754\ldots
\end{array}}
\]
These follow from explicit Clebsch--Gordan calculation for the
fundamental representation of $\SU(3)$ acting on the listed Casimir
levels.

\subsection{Scalar vs.\ channel-transfer bounds}

A naive polymer-resolvent bound replaces the off-diagonal matrix by
its $L^1$-norm (worst-case row sum):
\[
C_{\mathrm{step}}
\;:=\;
\sum_{i,j} A_{i\to j}
\;=\;
2\cdot(0.2083 + 0.0659 + 0.4612 + 0.1398)
\;=\;
0.875134.
\]
(The factor $2$ accounts for both directions.)

A sharper bound is the spectral radius of the channel transfer matrix.

\begin{definition}
The $4\times 4$ channel transfer matrix on $\mathcal L$ is
\[
T_{\mathrm{ch}}
\;:=\;
\begin{pmatrix}
0 & A_{0\to 1} & A_{0\to 2} & 0 \\
A_{0\to 1} & 0 & A_{1\to 2} & A_{1\to 3} \\
A_{0\to 2} & A_{1\to 2} & 0 & 0 \\
0 & A_{1\to 3} & 0 & 0
\end{pmatrix}.
\]
\end{definition}

\begin{theorem}[Channel spectral radius]\label{thm:channel-threshold}
The matrix $T_{\mathrm{ch}}$ has spectrum
\[
\spec(T_{\mathrm{ch}})
\;=\;
\{0.55016,\; -0.50450,\; -0.05158,\; 0.00592\},
\]
hence spectral radius $\rho_{\mathrm{ch}} = 0.55016\ldots$. The
sharpening factor over the scalar bound is
$C_{\mathrm{step}}/\rho_{\mathrm{ch}} = 1.5907$.
\end{theorem}

\begin{proof}
Direct diagonalization of the $4\times 4$ matrix. The dominant
eigenvalue is the Perron root; the negative subdominant eigenvalue
$-0.5045$ is close in magnitude, reflecting the bipartite-like
structure of the level coupling (alternating-sign symmetry of the
Casimir levels under $\hat\chi$-conjugation), but does not exceed
the Perron root in absolute value.
\end{proof}

\subsection{Polymer resolvent convergence}

\begin{corollary}[Polymer thresholds]\label{cor:polymer}
On a graph $\Gamma$ with adjacency-matrix spectral radius
$\rho(\mathrm{adj}(\Gamma))$, the polymer-resolvent series
\[
(M-V)^{-1} = \sum_{k=0}^\infty M^{-1}(VM^{-1})^k,
\qquad
M := \omega(\beta)\cdot\mathrm{Id},
\quad
V := T_{\mathrm{ch}}\otimes\mathrm{adj}(\Gamma),
\]
converges absolutely provided
\[
\omega(\beta)^{-1}\,\rho_{\mathrm{ch}}\,\rho(\mathrm{adj}(\Gamma)) < 1,
\]
i.e.\ provided
\[
\boxed{\;
\beta \;>\; \beta_{\mathrm{thresh}}(\Gamma)
\;:=\;
\frac{N}{2}\,\big(\rho_{\mathrm{ch}}\,\rho(\mathrm{adj}(\Gamma))\big)^2.
\;}
\]
\end{corollary}

\begin{proof}
The operator $VM^{-1}$ has norm bounded by
$\omega^{-1}\rho_{\mathrm{ch}}\rho(\mathrm{adj}(\Gamma))$. When this
is $<1$, the geometric series converges. The threshold $\beta$ is
obtained from $\omega(\beta) = \sqrt{\beta/N}$ via
$\omega^{-2}\rho_{\mathrm{ch}}^2\rho(\mathrm{adj}(\Gamma))^2 < 1$.
\end{proof}

\subsection{Explicit thresholds}

\begin{center}
\begin{tabular}{l c c c}
\toprule
Graph & $\rho(\mathrm{adj}(\Gamma))$ & Scalar
($C_{\mathrm{step}}$) & Channel ($\rho_{\mathrm{ch}}$) \\
\midrule
Chain (1D, deg 2) & $1.998$ & $\beta > 4.58$ & $\beta > 1.81$ \\
Cycle (1D, deg 2) & $2.000$ & $\beta > 4.60$ & $\beta > 1.82$ \\
2D grid (deg 4) & $3.759$ & $\beta > 16.23$ & $\beta > 6.42$ \\
4-regular & $4.000$ & $\beta > 18.38$ & $\beta > 7.27$ \\
6-regular & $6.000$ & $\beta > 41.36$ & $\beta > 16.35$ \\
\bottomrule
\end{tabular}
\end{center}

The headline result: on the 2D lattice (the simplest non-trivial
geometry), polymer-resolvent convergence is guaranteed for
$\beta \ge 6.42$. For the 4D lattice (random 4-regular as a
proxy for the unoriented adjacency), the threshold is $\beta \ge 7.27$.

\subsection{Deterministic exponential decay}

\begin{corollary}[Combes--Thomas type bound]\label{cor:CT}
For $\beta > \beta_{\mathrm{thresh}}(\Gamma)$, the lattice $H_1$ Green
function satisfies
\[
\|G^{(1)}_{bb'}(U)\|_{\mathrm{op}}
\;\le\;
\frac{1}{m_H^2(\beta)}\cdot\frac{1}{1-\tau}\cdot
e^{-m_0(\beta)\,d_\Gamma(b,b')},
\]
where $\tau = \omega^{-1}\rho_{\mathrm{ch}}\rho(\mathrm{adj}(\Gamma))
< 1$, $m_H^2(\beta) = \omega(\beta) = \sqrt{\beta/N}$, and the decay
rate satisfies
\[
m_0(\beta) \;\ge\; \log\Big(\omega(\beta)/
\big(\rho_{\mathrm{ch}}\rho(\mathrm{adj}(\Gamma))\big)\Big).
\]
\end{corollary}

\begin{proof}
Standard Combes--Thomas argument: insert the conjugation
$e^{m_0 d(\cdot,b)}\cdot e^{-m_0 d(\cdot,b)}$ into the resolvent
identity; the perturbation $e^{m_0 d}V e^{-m_0 d} - V$ is
controlled by $e^{m_0}$ on adjacency-step magnitude.
\end{proof}

This is the deterministic exponential-decay input needed by the
Helffer--Sj\"ostrand leakage framework (Section~\ref{sec:framework}).

\section{HK-Wilson calibration: independent verification}\label{sec:hk-cal}

The HK-Wilson calibration prescribes that the Wilson plaquette weight
$\exp(-\beta\phi)$ is comparable, in oscillation, to the heat-kernel
weight $p_{t(\beta)}$ at time
\[
t(\beta) \;=\; \frac{1}{\beta c_2} \;=\; \frac{N}{2\beta}.
\]
We verify this calibration independently via $\SU(2)$ boundary-layer
analysis.

\subsection{Exact $\SU(2)$ boundary-layer operator}

After eigenangle reduction, $\SU(2)$ at large $\beta$ reduces to the
half-line oscillator
\[
H_\eta \;=\; -\half\partial_u^2 + \big(\tfrac{1}{4}+2\eta\big)u^2,
\qquad u\in(0,\infty),\quad u\big|_{u=0}=0,
\]
with parameter $\eta\ge 0$ from the boundary-layer scaling. The exact
lowest gap is
\[
C_{\SU(2)}(\eta) \;=\; 2\sqrt{\tfrac{1}{2} + 4\eta}.
\]

\subsection{Convergence rate of the scaled transfer operator}

Defining the scaled operator $L_h := h^{-1}(I - S_h)$ where $S_h$ is
the heat-kernel transfer step at scale $h := \beta^{-1/2}$, finite
differences at $n=1600$ gridpoints and $\beta\in\{10^4, 10^5, 10^6\}$
give:
\[
\|L_h - H_\eta\|_{\mathrm{op}} \;\sim\; C\sqrt h,
\qquad
\frac{\|L_h - H_\eta\|}{\sqrt h} \to 0.40
\]
(empirical fit: $\|L_h-H_\eta\| = 0.428\,h^{0.51}$). The
predicted constant is
$\int_0^\infty t^2\Phi(-t)\,dt = \sqrt 2/(3\sqrt\pi) = 0.2660\ldots$,
and the empirical $0.40/0.27 \approx 1.5$ ratio is consistent with
multiplicative corrections from the angular projection.

\subsection{Exact spectral doubling under blocking}

A nontrivial cross-check: under $K \mapsto K^{2^m}$, the gap of any
non-degenerate transfer kernel should exactly double per step. Direct
computation at $\beta=5000$, $\eta=1$ gives:
\begin{center}
\begin{tabular}{cccl}
\toprule
$m$ & predicted gap & measured & rel.\ error \\
\midrule
0 & $5.787\times 10^{-2}$ & $5.787\times 10^{-2}$ & $0$ \\
1 & $1.157\times 10^{-1}$ & $1.157\times 10^{-1}$ & $3.6\times 10^{-14}$ \\
2 & $2.315\times 10^{-1}$ & $2.315\times 10^{-1}$ & $3.6\times 10^{-14}$ \\
3 & $4.630\times 10^{-1}$ & $4.630\times 10^{-1}$ & $4.8\times 10^{-14}$ \\
4 & $9.259\times 10^{-1}$ & $9.259\times 10^{-1}$ & $7.2\times 10^{-14}$ \\
\bottomrule
\end{tabular}
\end{center}
Agreement to machine precision confirms that the HK-Wilson calibration
$t(\beta) = N/(2\beta)$ is the unique scaling at which the blocking
flow is exact.

\section{Synthesis: place in the broader $\SU(N)$ mass-gap program}
\label{sec:framework}

\subsection{The $\sigma_{\mathrm{geom}}$ / Helffer--Sj\"ostrand
two-component decomposition}

The geometric-curvature framework states that the effective
coarse-grained Hessian satisfies
\[
\nabla^2 \bar H_{\mathrm{eff}}(\theta)
\;\succeq\;
\sigma_{\mathrm{geom}}\,I
\;-\;
\E_{\mu_\theta}\big[B^*(\mathcal A^{(1)})^{-1}B\big],
\]
where $\sigma_{\mathrm{geom}} = h^\vee/2 = N/2$ is the universal Lie
curvature (the Weyl/Jacobian curvature floor), $B = \nabla^2_{\theta\eta}S$
is the cross-Hessian, and $\mathcal A^{(1)}$ is the 1-form
(Helffer--Sj\"ostrand / Witten) operator on the internal variables.

\paragraph{First connection.}
The leading single-plaquette gap $\sqrt{2\beta/N}$ of
Proposition~\ref{prop:sun-pattern} is precisely the dynamical
contribution to $\sigma_{\mathrm{geom}}$ at large $\beta$ in the radial
sector. The static $\sigma_{\mathrm{geom}} = N/2$ is the
$\beta\to\infty$ limit of the radial gap divided by $\sqrt\beta$
(squared) and rescaled. Theorem~\ref{thm:su3-asym} therefore
provides the explicit large-$\beta$ expansion of the curvature floor
to three orders.

\paragraph{Second connection.}
The polymer threshold $\beta_{\mathrm{thresh}}(\Gamma)$ of
Theorem~\ref{thm:channel-threshold} gives the explicit condition
under which the 1-form resolvent $\mathcal A^{(1)}$ admits a
deterministic exponential-decay bound (Corollary~\ref{cor:CT}).
Above the threshold, the Schur leakage term
$B^*(\mathcal A^{(1)})^{-1}B$ is itself locally bounded, with
exponentially-decaying off-diagonal entries.

\subsection{The intermediate coupling window}

Combining the two anchors gives the following picture for $\SU(3)$
in $d = 4$:
\begin{itemize}
\item\textbf{Small $\beta$:} the Dobrushin LSI of
\cite{ShenZhuZhu2023} and the refinement of
\cite{CaoNissimSheffield2025} establish a volume-uniform LSI for
$\beta < \beta_{\mathrm{Dob}} \approx 0.175$ (Dobrushin matrix
condition).

\item\textbf{Large $\beta$:} Corollary~\ref{cor:polymer} establishes
deterministic resolvent decay for
$\beta > \beta_{\mathrm{thresh}} \approx 7.3$ (on the 4-regular
lattice graph).

\item\textbf{Intermediate window:} $0.175 < \beta < 7.3$ remains open.
\end{itemize}

This intermediate window is significantly smaller than in prior
literature, where the corresponding scalar bound gives
$\beta_{\mathrm{thresh}}^{\mathrm{scalar}} \approx 18$. The channel
refinement closes a factor of $\sim 2.5$.

\subsection{Asymptotic freedom trajectory}

Under the asymptotic-freedom flow
$\beta(a) \sim 1/(b_0\log(1/a\Lambda))$ with $b_0 = 11N/(48\pi^2)$,
the polymer threshold condition $\beta(a) > \beta_{\mathrm{thresh}}$
is satisfied at all sufficiently small $a$, since
$\beta(a)\to\infty$ as $a\to 0$. Therefore the polymer-resolvent
decay is \emph{automatic} in the continuum limit; the
intermediate-coupling obstruction is exclusively a finite-$\beta$
phenomenon. The remaining work is the Mosco / form-convergence step
relating the lattice gap to its continuum limit:
\[
m_{\mathrm{phys}}(a) \;=\; \frac{\Gap(a)}{a}
\;\sim\;
\frac{1}{a}\sqrt{\frac{2\beta(a)}{N}}
\;\sim\;
\frac{1}{a\sqrt{\log(1/a\Lambda)}}.
\]
This naively diverges as $a\to 0$, requiring an appropriate
renormalization-group rescaling to extract the physical gap.

\section{Open directions}

\begin{enumerate}[label=(O\arabic*)]
\item\textbf{Closed-form $r_N$.} Identify the constant
$O(\beta^0)$ term in $\Gap_{\SU(N)}$ as an explicit
group-theoretic combination. Conjecture: $r_N$ admits a closed form
involving the dimension $d_G = N^2-1$ and the $d$-symbol norm
$\sum_{abc}d_{abc}^2$.

\item\textbf{Closed-form $c_1$ for general $N$.} Extend the
$311\sqrt 6/9216$ derivation. The structural ingredients (Wilson
Taylor coefficients, angular averages, Laguerre matrix elements in
$d_G$ dimensions) are identical; the bottleneck is the explicit
$\E_{S^{d_G-1}}[(\Tr A^3)^2]$ for $N\ge 4$.

\item\textbf{Closing the intermediate coupling window.}
$\beta \in (0.175, 7.3)$ for $\SU(3)$ in $d=4$ remains the only
unresolved finite-$\beta$ window. The $\sigma_{\mathrm{geom}}$ floor
plus a typicality estimate on the small-field event is the natural
attack.

\item\textbf{Channel threshold sharpening.} Extending
$T_{\mathrm{ch}}$ from $4\times 4$ to $8\times 8$ or larger should
further reduce $\rho_{\mathrm{ch}}$, sharpening the threshold further.

\item\textbf{Mosco convergence on the AF trajectory.} Combined with
the polymer decay, the universal radial gap of Theorem~\ref{thm:su3-asym}
gives the input for a uniform-in-$\Lambda$ Poincar\'e inequality. The
Mosco / form-convergence step closing the continuum limit remains the
final piece.
\end{enumerate}

\appendix

\section{Explicit pentadiagonal entries of $\hat T^2$}\label{app:T2}

We verify the entries of $\hat T^2$ used in
Section~\ref{ssec:c1}(b) directly from the tridiagonal Jacobi matrix
of Lemma~\ref{lem:laguerre-8}. The nonzero entries of $\hat T$ are
\[
\hat T_{n,n} \;=\; 2n+4,
\qquad
\hat T_{n,n+1} \;=\; \hat T_{n+1,n} \;=\; -\sqrt{(n+1)(n+4)}.
\]

\paragraph{Pentadiagonal structure.}
Since $\hat T$ is tridiagonal, $(\hat T^2)_{mn} = \sum_k \hat T_{mk}\hat T_{kn}$
vanishes unless $|m-n|\le 2$. The three nontrivial families are computed
from~\eqref{eq:T2-entries}:
\begin{align*}
\hat T^2_{n,n}
&\;=\; \hat T_{n,n-1}^2 + \hat T_{n,n}^2 + \hat T_{n,n+1}^2 \\
&\;=\; n(n+3) + (2n+4)^2 + (n+1)(n+4) \\
&\;=\; 6n^2 + 24n + 20, \\[3pt]
\hat T^2_{n+1,n}
&\;=\; \hat T_{n+1,n}\,\hat T_{n,n} + \hat T_{n+1,n+1}\,\hat T_{n+1,n} \\
&\;=\; \hat T_{n+1,n}\bigl(\hat T_{n,n}+\hat T_{n+1,n+1}\bigr) \\
&\;=\; -\sqrt{(n+1)(n+4)}\,(4n+10),\\[3pt]
\hat T^2_{n+2,n}
&\;=\; \hat T_{n+2,n+1}\,\hat T_{n+1,n}
\;=\; \sqrt{(n+1)(n+2)(n+4)(n+5)}.
\end{align*}

\paragraph{Numerical values used in $S_0$, $S_1$.}
\begin{center}
\begin{tabular}{c|rrrrr}
\toprule
$n$ & $0$ & $1$ & $2$ & $3$ & $4$ \\
\midrule
$\hat T^2_{n,n}$
& $20$ & $50$ & $92$ & $146$ & $212$ \\
$\hat T^2_{n+1,n}$
& $-20$ & $-14\sqrt{10}$ & $-54\sqrt{2}$ & $-44\sqrt{7}$ & $-52\sqrt{10}$ \\
$\hat T^2_{n+2,n}$
& $2\sqrt{10}$ & $6\sqrt{5}$ & $6\sqrt{14}$ & $4\sqrt{70}$ & $12\sqrt{15}$ \\
\bottomrule
\end{tabular}
\end{center}
(Squared values appearing in the Rayleigh--Schr\"odinger sums:
$|\hat T^2_{1,0}|^2 = 400$,
$|\hat T^2_{2,0}|^2 = 40$,
$|\hat T^2_{2,1}|^2 = 1960$,
$|\hat T^2_{3,1}|^2 = 180$.)

\paragraph{Verification of $S_0, S_1$.}
\begin{align*}
S_0 &\;=\; \frac{|\hat T^2_{1,0}|^2}{0-1} + \frac{|\hat T^2_{2,0}|^2}{0-2}
\;=\; \frac{400}{-1} + \frac{40}{-2} \;=\; -420,\\
S_1 &\;=\; \frac{|\hat T^2_{0,1}|^2}{1-0} + \frac{|\hat T^2_{2,1}|^2}{1-2} + \frac{|\hat T^2_{3,1}|^2}{1-3}\\
&\;=\; \frac{400}{1} + \frac{1960}{-1} + \frac{180}{-2}
\;=\; 400 - 1960 - 90 \;=\; -1650.
\end{align*}
Hence $S_1 - S_0 = -1230$, giving~\eqref{eq:c1-quart}.

\paragraph{Reproducibility.}
The full computation reduces to integer arithmetic (after factoring out
the common $\sqrt{6}$). A short \texttt{sympy} script that reproduces
$c_1 = -311\sqrt{6}/9216$ from these entries is given below.

\begin{verbatim}
from sympy import sqrt, Rational, simplify, Integer
# Pentadiagonal T^2 entries (squared, used in sums)
T2sq = {(1,0): Integer(400), (2,0): Integer(40),
        (2,1): Integer(1960), (3,1): Integer(180)}
# S_n = sum_{m != n} (T^2_{mn})^2 / (n - m)
S0 = sum(v / Integer(0 - m) for (m, n), v in T2sq.items() if n == 0)
S1 = sum(v / Integer(1 - m) for (m, n), v in T2sq.items() if n == 1)
S1 += T2sq[(1, 0)] / Integer(1)            # T^2_{0,1} = T^2_{1,0}
# c_1^{qrt-2nd} = (S1 - S0) * sqrt(6) / 18432
c1_quart = simplify((S1 - S0) * sqrt(6) / 18432)
c1_sext  = Rational(19, 576) * sqrt(6)     # = 304*sqrt(6)/9216
c1 = simplify(c1_sext + c1_quart)
print(c1)   # -311*sqrt(6)/9216
\end{verbatim}

\section*{Acknowledgments}

This paper synthesizes computational results from notebooks
GOODRESULTS\{1,\ldots,5\} (February 2026), in conjunction with the
broader $\sigma_{\mathrm{geom}}$ / Helffer--Sj\"ostrand framework
developed in companion notes.

\begin{thebibliography}{99}

\bibitem{ShenZhuZhu2023} H.\ Shen, R.\ Zhu, X.\ Zhu (2023).
\emph{Stochastic quantization of Yang--Mills--Higgs on the 2-d torus.}
arXiv:2308.00040.

\bibitem{CaoNissimSheffield2025} S.\ Cao, A.\ D.\ Nissim, S.\ Sheffield
(2025). \emph{Continuum Yang--Mills mass gap from the lattice via a
log-Sobolev inequality.} (Preprint.)

\bibitem{notebooks} GOODRESULTS\{1,\ldots,5\}.ipynb. Internal
computational notebooks, February 2026.

\bibitem{Tlas2026} T.\ Tlas (2026). \emph{Gaussian concentration for
$\SU(N)$ Yang--Mills.} arXiv:2603.25236.

\bibitem{HelfferSjostrand} B.\ Helffer, J.\ Sj\"ostrand (1994).
\emph{Decay of correlations for unbounded interacting spin systems.}

\bibitem{BBD} R.\ Bauerschmidt, T.\ Bodineau, B.\ Dagallier (2025).
\emph{Stochastic quantization, log-Sobolev inequalities and Yang--Mills.}
arXiv:2503.24372.

\end{thebibliography}

\end{document}
